\chapter{Layout: grid and absolute positioning}
\label{ch:grid}

\begin{aside}
Flexbox handles 1-D layouts well. For 2-D, grid is the
answer. For ``put this exactly here'' moments, absolute
positioning escapes layout entirely.
\end{aside}

\section{Grid basics}

\begin{lstlisting}
auto layout = grid({
    .rows = "auto 1fr auto",
    .cols = "200 1fr"
}, {
    {0, 0, header() | span(2)},
    {1, 0, sidebar()},
    {1, 1, main_content()},
    {2, 0, footer() | span(2)},
});
\end{lstlisting}

Each child occupies cells in the grid; \code{span} extends
across rows or columns.

\section{Track sizing}

\begin{itemize}[leftmargin=*]
  \item Fixed: literal cell count.
  \item \code{auto}: shrink to content.
  \item \code{1fr}: fractional, share leftover.
  \item \code{minmax(a, b)}: between min and max.
\end{itemize}

\section{Areas}

A named-area syntax:

\begin{lstlisting}
auto layout = grid({
    .areas = R"(
        header header
        side   main
        footer footer
    )"
}, {
    {.area = "header", .child = header()},
    // ...
});
\end{lstlisting}

Readable; reflects the layout shape.

\section{Absolute positioning}

For overlays, tooltips, popovers --- escape layout:

\begin{lstlisting}
auto popup = absolute(card, .x = 20, .y = 5);
\end{lstlisting}

The popup paints at the given coordinates regardless of
parent.

\section{Relative positioning}

\code{relative} shifts an element from where layout placed
it. Useful for animations or decorative offsets.

\section{Z-index}

For overlays with stacking order:

\begin{lstlisting}
auto over = card | z(100);
\end{lstlisting}

Higher z paints later (on top). See Chapter~\ref{ch:zorder}.

\section{Recap}

\begin{recap}
\begin{itemize}[leftmargin=*]
  \item Grid for 2-D; flexbox for 1-D.
  \item Track sizes: fixed, auto, fr, minmax.
  \item Named areas for readable layouts.
  \item Absolute for overlays and popovers.
  \item Z-index controls stacking.
\end{itemize}
\end{recap}

\section*{Workshop}

\begin{enumerate}[leftmargin=*]
  \item Build a holy-grail layout with grid.
  \item Use named areas.
  \item Add an absolutely-positioned tooltip.
\end{enumerate}
